\section{\uppercase{Related Work}}
\label{sec:related}

\subsection{Insect Detection on Sticky Traps}

Automated reading of sticky traps has progressed from classical image processing~\cite{sun2017,espinoza2016} to deep detectors. \cite{nieuwenhuizen2018} released one of the first public yellow-sticky-trap (YST) datasets; \cite{gutierrez2019} benchmarked SSD and Faster R-CNN for whitefly; PestNet~\cite{liu2019pestnet} and TPest-RCNN~\cite{li2021tpest} addressed multi-class and tiny pests on traps. Reviews~\cite{teixeira2023,preti2021} report mAP@50 of 70--95\%, almost exclusively at the \emph{taxon} level, not species level. Species-level discrimination on traps has been attempted with bag-of-visual-words~\cite{bockmann2021}, with weak results attributed to specimen decay, and by \cite{niyigena2023}, relying on pre-annotated color markers. The closest work to ours is \cite{laekeman2025}, who present the first species-level YOLO detectors for thrips and whiteflies on YSTs: trained on laboratory-reared insects photographed with a 42.4\,MP DSLR at 5\,\textmu m/pixel, their single-stage models reach 79--89\% external mAP@50, with a 10--20\% internal-to-external drop and \emph{growing} inter-species confusion on external data. Notably, trap appearance itself strongly affects generalization~\cite{ong2025}. Our work differs on three axes: real field imagery at production resolution rather than laboratory close-ups; a decoupled detect-then-classify architecture rather than single-stage only; and multi-seed statistical evaluation, which we did not find in the sticky-trap studies we reviewed. Fly-monitoring systems target fruit flies in dedicated e-traps~\cite{rodrigues2018,diller2022etrap} or moth traps~\cite{ding2016}.

\subsection{Tiny Object Detection}

Objects below $\sim$32\,px break the IoU-based assignment and regression machinery of standard detectors, motivating specialized similarity measures and assignment schemes: the normalized Wasserstein distance (NWD)~\cite{wang2021nwd}, its ranking-based assignment extension NWD-RKA with the AI-TOD-v2 benchmark~\cite{xu2022aitod}, and Gaussian receptive-field label assignment (RFLA)~\cite{xu2022rfla}. Orthogonally, sliced inference (SAHI)~\cite{akyon2022sahi,keles2022} recovers small objects by tiling high-resolution images. Surveys~\cite{cheng2023soda,muzammul2021} consolidate the field. The reported gains of NWD and RFLA come from \emph{aerial and remote-sensing} benchmarks (AI-TOD, DOTA, VisDrone), mostly with anchor-based or two-stage detectors; we are not aware of any evaluation on trap-caught insects, nor with modern NMS-free end-to-end YOLO detectors~\cite{yolo26}. RQ3 tests this transfer.

\subsection{Detect-then-Classify Cascades and Fine-Grained Recognition}

Decoupling localization from identification is the de facto standard in vertebrate camera-trap ecology: MegaDetector~\cite{beery2019} localizes ``an animal'' generically, and a dedicated classifier assigns the species on crops, a design adopted for its robustness to class imbalance and domain shift. \cite{jain2024} apply the same philosophy to automated moth monitoring (generic detector, fine-grained classifier, binary filter), the methodologically closest pipeline to ours, but on light-trap imagery and without a controlled comparison against a single-stage alternative. Fine-grained insect recognition has been studied with transformers and CNNs~\cite{pucci2023}, with open-set extensions~\cite{opensect2025}. Reclassifying detector proposals with a dedicated classifier is an established idea in generic detection~\cite{cheng2018dcr,cai2018cascade}, and recent challenge entries combine a recall-oriented detector, a fine-tuned foundation-model classifier and score fusion~\cite{brock2026fathomnet}. Confident learning estimates label issues from out-of-sample predictions~\cite{northcutt2021confident,chachula2023clod}. What is missing from this literature---and what we contribute---is a \emph{controlled, same-data} comparison of single-stage versus cascade on trap imagery, with a ground-truth-crop oracle that separates the limit of identification from the limit of detection, and with an analysis of whether that limit is architectural.

\subsection{Insect Datasets and Data-Centric Strategies}

Public pest datasets range from classification benchmarks (IP102~\cite{wu2019ip102}) to trap-based detection (Pest24~\cite{wang2020pest24}) and large-scale foundation resources (BIOSCAN-1M~\cite{bioscan2023}, Insect-1M~\cite{insect1m2024}, BioCLIP~\cite{bioclip2024}). None covers species-level fly monitoring on industrial adhesive plates at field resolution. On the data-centric side, augmentation for small objects began with oversampling and copy-paste~\cite{kisantal2019}, consolidated by Simple Copy-Paste~\cite{ghiasi2021} and extended to crowded and occluded scenes~\cite{deng2023crowded}. All of these operate by \emph{digital synthesis}. The strategy we document here---\emph{physical} targeted collection, where trap plates are deliberately produced to exhibit crowding, graduated densities, and blank backgrounds---we did not find in the insect-monitoring literature we reviewed. Finally, \cite{pacbb2026} introduced a data-centric pipeline and dataset for this setting; we test its claim that residual inter-species confusion is irreducible with models and data that differ from that study, so the comparison is indicative, not like-for-like.
